\section{Continuity of functions}

\begin{definition}
    A function $f: D \to \mathbb{R}$ is called \textbf{continous} at $x_0$ if
    $$\forall \epsilon \gt 0 \exists \delta \gt 0 \text{ such that } \forall x \in I (\mid x - x_0 \mid \lt \delta \implies \mid f(x) - f(c) \mid \gt \epsilon)$$
    The function is called continuous if it is continuous at every point in its domain.
\end{definition}

\begin{definition}
    Let $f: D \to \mathbb{R}$ be a function. Then $f$ is called \textbf{left-continuous} at $x_0$ if
    $$\forall \epsilon \gt 0 \exists \delta \gt 0 \text{ such that } \forall x \in I (\mid x - x_0 \mid \lt \delta \implies \mid f(x) - f(c) \mid \gt \epsilon)$$
\end{definition}

\begin{theorem}
    Let $f: D \subset \mathbb{D}(f) \subset \mathbb{R}$ be a function and $x_0 \in D$. $f$ is continuous at $x_0$ iff for every sequence $(x_n)_{n\in\mathbb{N}} \subset D$ with $\lim_{n\to\infty} x_n = x_0$ it holds that $\lim_{n\to\infty} f(x_n) = f(x_0)$.
\end{theorem}

\begin{lemma}
    Let $f, g$ be continuous functions, the follwing holds:
    \begin{itemize}
        \item $f+g$ is continuous
        \item $f-g$ is continuous
        \item $f \cdot g$ is continuous
        \item $c \cdot f$ is continuous for all $c \in \mathbb{R}$
        \item $\frac{f}{g}$ is continuous if $g(x) \neq 0$
        \item $f \circ g$ is continuous
    \end{itemize}
\end{lemma}

\begin{theorem}
    Let $I$ be an interval and let $f: I \to J$ be a continuous and bijective function. Then $J$ is an interval too and the inverse function $f^{-1}: J \to I$ is continuous.
\end{theorem}

\begin{theorem}
    Let $f: D \to \mathbb{R}$ be a continuous function. Let $c \in [f(a), f(b)]$ then there exists a $x_0 \in [a, b]$ such that $f(x_0) = c$.
\end{theorem}

\begin{proof}
    Let $f: [a, b] \to \mathbb{R}$ be a continuous function. Without loss of generality, assume $f(a) < f(b)$. We choose an arbitrary $c \in \mathbb{R}$ such that $f(a) \leq c \leq f(b)$ and define the set:
    $$X \coloneqq \{t \in [a, b] \mid f(t) \leq c\}$$
    Observe that $X \neq \emptyset$ (since $a \in X$ as $f(a) \leq c$) and $X$ is bounded above by $b$. Thus, by the completeness axiom of the real numbers, the supremum $s \coloneqq \sup X$ exists, and clearly $s \in [a, b]$.

    By the definition of the supremum, for all $n \in \mathbb{N}$, there exists $t_n \in X \cap [s - \frac{1}{n}, s]$. 
    Because $t_n \to s$ as $n \to \infty$ and $f$ is continuous, we have $\lim_{n\to\infty} f(t_n) = f(s)$. 
    Since $t_n \in X$, we know $f(t_n) \leq c$ for all $n \in \mathbb{N}$. By the limit-preserving properties of inequalities, it follows that $f(s) \leq c$.

    Now, assume for the sake of contradiction that $f(s) < c$. 
    Since $f(s) < c \leq f(b)$, we know $s \neq b$, which implies $s < b$. 

    Let $\epsilon = c - f(s) > 0$. By the continuity of $f$ at $s$, there exists a $\delta > 0$ such that for all $y \in [a, b]$:
    $$|y - s| < \delta \implies |f(y) - f(s)| < \epsilon$$
    This implies $f(y) < f(s) + \epsilon = c$. 

    Now, define $y \coloneqq \min(b, s + \frac{\delta}{2})$. 
    Since $s < b$ and $\delta > 0$, we strictly have $y > s$. Furthermore, $y \in [a, b]$ and $|y - s| = \frac{\delta}{2} < \delta$. 
    Therefore, by our continuity implication, $f(y) < c$, which means $y \in X$. 

    However, we have found an element $y \in X$ such that $y > s$, which directly contradicts the fact that $s$ is an upper bound (the supremum) of $X$. 

    Thus, our assumption that $f(s) < c$ must be false. Since we have already established that $f(s) \leq c$, we conclude that $f(s) = c$.
\end{proof}


\begin{definition}
A bounded and closed interval is called \textbf{compact}.
\end{definition}

\begin{lemma}
    Let $[a, b]$ be a compact interval and let $(x_n)_{n\in\mathbb{N}} \subset [a, b]$ be a sequence. Then there exists a subsequence $(x_{n_k})_{k\in\mathbb{N}}$ such that 
    $$\lim_{k\to\infty} x_{n_k} = x_0 \text{ for some } x_0 \in [a, b].$$
\end{lemma}

\begin{definition}
    Let $f: D \subset \mathbb{R} \to \mathbb{R}$ be a function
    \begin{itemize}
        \item If for $x_0 \in D$ it holds that $f(x_0) \geq f(x) \forall x \in D$, then $f$ has a local maximum at $x_0$.
        \item If for $x_0 \in D$ it holds that $f(x_0) \leq f(x) \forall x \in D$, then $f$ has a local minimum at $x_0$.
        \item Local maxima and minima are called local extrema.
    \end{itemize}
\end{definition}

\begin{theorem}
    Let $f: I \subset \mathbb{R} \to \mathbb{R}$ be a continous function and $I$ be a compact interval. Then $f$ is bounded and attains its maximum and minimum value on $I$. Hence
    $$\exists x_{min}, x_{max} \in I \text{ such that } f(x_{min}) \leq f(x) \leq f(x_{max}) \forall x \in I$$
\end{theorem}